\documentclass[book.tex]{subfiles}

\begin{document}

\chapter{Noether Networks}
\label{chap:noether_networks}
\noindent\rule{11cm}{0.4pt} \\
\textbf{Prerequisites:} \autoref{chap:pde}, \autoref{chap:simple_physics} \\
\textbf{Difficulty Level:} ** \\
\noindent\rule{11cm}{0.4pt}
\vspace{5mm}

%"""
%From CNNs to attention mechanisms, encoding inductive biases into neural networks has been a fruitful source of improvement in machine learning. Adding auxiliary losses to the main objective function is a general way of encoding biases that can help networks learn better representations. However, since auxiliary losses are minimized only on training data, they suffer from the same generalization gap as regular task losses. Moreover, by adding a term to the loss function, the model optimizes a different objective than the one we care about. To fix both problems we propose tailoring: a framework that fine-tunes the model for each particular input with unsupervised objectives within the prediction function, both at training and test time. The flexibility of the framework allows a variety of applications, from leveraging contrastive losses, to softly encoding physics-based inductive biases or improving adversarial defenses. In the second part of the talk, we’ll discuss how we can discover inductive biases for sequential prediction problems. Inspired by Noether’s theorem, equating symmetries and conservation laws, we propose to meta-learn useful conservations imposed using the tailoring framework. We show, theoretically and experimentally, that Noether Networks improve prediction quality, providing a general framework for discovering inductive biases in sequential problems.
%"""

Noether's theorem draws a deep link between symmetries of a system and conservation laws. Recent work suggests that Noether's theorem provides a framework to help learn useful conservation laws for machine learning applications. This provides a potentially general methodology to automatically generate inductive biases for a system \cite{alet2020tailoring}, \cite{alet2021noether}


\section{A Brief Review of Noether's Theorem}

Noether's theorem is one of the most foundational concepts in modern physics. Roughly stated, it says that to every symmetry of a physical system, there is a paired conservation law. For example, if translation symmetry holds, then linear momentum is conserved. 


\end{document}